\documentclass[11pt,a4paper]{article}

\usepackage[margin=1.05in]{geometry}
\usepackage{amsmath,amssymb,amsthm}
\usepackage{mathtools}
\usepackage{stmaryrd}
\usepackage{booktabs}
\usepackage{array}
\usepackage{xcolor}
\usepackage{listings}
\usepackage{tikz}
\usetikzlibrary{arrows.meta,positioning,calc}
\usepackage[colorlinks=true,linkcolor=black,citecolor=black,urlcolor=blue]{hyperref}
\usepackage{enumitem}
\usepackage[protrusion=true,expansion=false]{microtype}
\setlength{\emergencystretch}{2.5em}

\theoremstyle{definition}
\newtheorem{requirement}{Requirement}[section]
\newtheorem{decision}[requirement]{Decision}
\newtheorem{definition}[requirement]{Definition}
\theoremstyle{remark}
\newtheorem{remark}[requirement]{Remark}
\newtheorem{finding}[requirement]{Finding}

% NOTATION RULES: the rho calculus is never written with the Greek letter;
% quoting is @P, dropping is *x.
\newcommand{\rhoc}{\textsf{rho}}
\newcommand{\Rnn}{\mathbb{R}_{\geq 0}}
\newcommand{\Qnn}{\mathbb{Q}_{\geq 0}}
\newcommand{\Pit}{\Pi}
\newcommand{\Dur}{\Delta}
\newcommand{\nm}[3]{\langle @#1,\,#2,\,#3\rangle}
\newcommand{\crisp}[1]{\lceil #1 \rceil}
\newcommand{\len}[1]{|#1|}
\newcommand{\nulld}{\epsilon}
\newcommand{\rest}{\mathsf{r}}
\newcommand{\MUST}{\textsc{must}}
\newcommand{\SHOULD}{\textsc{should}}
\newcommand{\MAY}{\textsc{may}}
\newcommand{\cmd}[1]{\texttt{#1}}
\newcommand{\atm}[1]{\mathsf{#1}}
\newcommand{\V}{\mathbb{V}}
\newcommand{\Rec}{\mathsf{Rec}}
\newcommand{\Fan}{\mathsf{Fan}}
\newcommand{\Server}{\mathsf{Server}}
\newcommand{\Voice}{\mathsf{Voice}}
\newcommand{\Call}{\mathsf{Call}}
\newcommand{\Resp}{\mathsf{Resp}}

\definecolor{codebg}{gray}{0.96}
\lstdefinelanguage{score}{
  morekeywords={score,pitches,durations,timbres,chance,def,factor,machine,table,
    for,where,par,in,at,play,line,import,proc,name,timbre,pitch,dur,clause,
    true,false,nu,mu,held,carry,prev,step,back,last,passes,leads},
  sensitive=true,
  morecomment=[l]{//},
  morestring=[b]",
}
\lstdefinelanguage{Rust}{
  morekeywords={pub,struct,enum,impl,fn,trait,let,mut,use,for,where,self,Self,Box,Vec,Result},
  sensitive=true, morecomment=[l]{//}, morestring=[b]",
}
\lstset{
  basicstyle=\ttfamily\small,
  backgroundcolor=\color{codebg},
  frame=none,
  keywordstyle=\bfseries,
  commentstyle=\itshape\color{black!60},
  columns=fullflexible,
  keepspaces=true,
  showstringspaces=false,
  xleftmargin=6pt,
  aboveskip=6pt, belowskip=6pt,
  breaklines=true,
}

\title{\bfseries The f1r3score Player\\[4pt]
\large Specification for a Rust tool that reads a textual score\\
and plays a run of it}

\author{F1R3FLY.io}

\date{Specification --- 25 September 2026\\[2pt]
\normalsize Revision 1. Implements \emph{Scores as Processes} (second version),
\texttt{publications/f1r3score}}

\begin{document}
\maketitle

\begin{abstract}
This specification defines \cmd{f1r3score}, a Rust command-line player and
library for the score calculus of \emph{Scores as Processes}. The player reads a
score written in a textual surface syntax, elaborates it to a closed term of the
calculus, and executes one run: it enumerates candidates at each location,
evaluates the graded where-clauses on them, groups them into contention sets,
selects a maximal matching with probability proportional to the product of its
weights, and stamps the emitted notes with onsets computed from the data. Its
outputs are the performance (the multiset of timed notes of positive length), a
replayable trace, a Standard MIDI File, and optionally live MIDI.

Three properties organise the design. \emph{Faithfulness}: the engine implements
the reduction and resolution rules of the note literally, with no knowledge of
voices, fans or idioms; those exist only as library definitions in the surface
syntax. \emph{Exactness}: weights are exact rationals and chance is consumed by
exact interval decoding, so a score, a chance source and a scheduler determine
a run bit for bit on every platform. \emph{Conformance}: the Python reference
simulator \texttt{skeinsim.py} and its thirteen tests are ported as the
acceptance suite, with the same statistics and the same thresholds.

The document fixes the textual format (\S\ref{sec:syntax}), the crate
architecture and data model (\S\ref{sec:arch}), the clause language and its
checkable fragment (\S\ref{sec:clauses}), the engine (\S\ref{sec:engine}), the
chance sources (\S\ref{sec:chance}), the outputs and command line
(\S\S\ref{sec:outputs}--\ref{sec:cli}), conformance (\S\ref{sec:conformance}),
a phased plan (\S\ref{sec:phases}) and the decisions still open
(\S\ref{sec:open}).
\end{abstract}

\tableofcontents

%==============================================================================
\section{Purpose and scope}\label{sec:scope}

\subsection{What the tool does}

Given a file \texttt{piece.score}, the command
\begin{lstlisting}
f1r3score play piece.score --chance prng:7 --notes 64 --midi piece.mid
\end{lstlisting}
parses and elaborates the score, runs it until 64 notes of positive length have
sounded or nothing more can fire, writes the performance to standard output,
and writes a MIDI file. The same run, replayed from its trace, reproduces the
same performance exactly.

\subsection{What was reviewed}

\begin{itemize}[leftmargin=*,itemsep=2pt]
\item \texttt{publications/f1r3score/score-calculus.tex} (second version): the
  normative semantics. Section and result numbers below refer to it.
\item \texttt{publications/f1r3score/skeinsim.py}: the stdlib-only reference
  simulator (716 lines, tests T0--T12). Its data model --- a soup indexed by
  (quote, timbre), hash-consed quotes, generic candidate generation --- is the
  model this specification adopts.
\item \texttt{publications/fuzzyware/graded-where-pln.tex} (Paper~II): the
  two formula sorts, per-candidate evaluation, contention sets, and the
  resolution algebra $(\Rnn,+,\times,0,1)$.
\item \texttt{publications/f1r3skein/f1r3skein-drawn-distributions.tex}: the
  weighted state machines drawn on the instrument, their canonical form, the
  \texttt{Machines} Theory, and exact interval decoding (Algorithm~1).
\item \texttt{F1R3Games/F1R3Skein}: the Rust workspace (\texttt{spigot\_stream},
  \texttt{dual\_spigot}, \texttt{spigot\_midi}, \texttt{leap\_spigot}), for
  reusable digit streams and MIDI code.
\item The rho-weighted crate and its design (\texttt{publications/rho-life/Gillespie}),
  for the ideal-versus-checkable discipline expressed as types.
\end{itemize}

\subsection{What the tool is not}

\begin{itemize}[leftmargin=*,itemsep=2pt]
\item \textbf{Not an interpreter on a shard.} Following the precedent set for
  the weighted-GSLT simulator, the player is a separate, offline execution path.
  It has no dependency on \texttt{f1r3node-rust}, RSpace or consensus.
\item \textbf{Not F1R3Skein.} The score calculus generalises an idea from the
  instrument; it is not a formalisation of it. The instrument's streams are one
  score among many (Remark~4.6 of the note), and the player treats them so: a
  drawn machine can be imported as a clause factor, nothing more.
\item \textbf{Out of scope:} tempo as a graded dimension, dynamics and
  articulation, timbre change in a communication, complex grading, and learning
  styles from a corpus --- all listed as open threads in the note. The renderer
  takes a fixed tempo as a rendering parameter only.
\end{itemize}

\subsection{Conventions}

\MUST, \SHOULD{} and \MAY{} have their RFC~2119 meanings. Requirements are
numbered R$n$ and decisions D$n$. Rust identifiers are indicative; the phase
exit criteria are normative.

%==============================================================================
\section{The semantic contract}\label{sec:contract}

This section restates, as obligations on the implementation, what the note
defines. Anything not stated here follows the note.

\begin{requirement}[Data]\label{req:data}
A score fixes a finite pitch alphabet $\Pit$ containing the rest $\rest$, a
finite duration alphabet $\Dur$ with lengths $\len{\cdot}:\Dur\to\Qnn$ in whole
notes, containing the null duration $\nulld$ with $\len{\nulld}=0$, and a finite
set of timbres. $\Pit$ and $\Dur$ \MUST{} be disjoint, since a datum fixes the
polarity of its name.
\end{requirement}

\begin{requirement}[Terms]\label{req:terms}
Names are variables or triples $\nm{P}{\tau}{d}$; processes are $0$, receipts
$\mathsf{for}(y\leftarrow x\;\mathsf{where}\;\psi)P$, joins (chord receipts,
Definition~5.5), messages $x!(Q,\tau,d)$, parallel composition and
$\ast x$. There is no replication, no \textsf{new}, and no recursion
former.
\end{requirement}

\begin{requirement}[Congruence]\label{req:cong}
Structural congruence is the least congruence closed under
$\alpha$-conversion making $(\mid,0)$ a commutative monoid and satisfying
$\ast\nm{P}{\tau}{d}\equiv P$. Two names are equivalent when their quoted
processes are congruent and their decorations are equal. The implementation
\MUST{} decide name equivalence exactly.
\end{requirement}

\begin{requirement}[Communication]\label{req:comm}
A candidate is a top-level receipt and a top-level message at the same location
and timbre, of opposite polarities, whose payload timbre equals the subject
timbre (\textsc{Comm}$_1$/\textsc{Comm}$_2$). Firing it emits the note
$(\tau,\pi,\delta)$ --- the pitch from whichever side carries one --- and
continues with $P\{\nm{Q}{\tau}{d}/y\}$. Every reduction emits exactly one note;
there are no silent steps (Proposition~2.4).
\end{requirement}

\begin{requirement}[Weights and resolution]\label{req:resolve}
A candidate at a receipt with clause $\psi$ is \emph{available} iff
$\psi\llbracket c\rrbracket>0$. Contention sets are the connected components of
the relation ``shares a receipt or a message'' over available candidates. A set
is resolved by choosing a maximal matching $\mathcal M$ with probability
proportional to $\prod_{c\in\mathcal M}\psi_c\llbracket c\rrbracket$ and firing
all of it (Definition~3.3).
\end{requirement}

\begin{requirement}[Time]\label{req:time}
Every top-level receipt and message carries a stamp in $\Qnn$. A communication
has onset $\max(t_r,t_m)$, emits the timed note $(t,\tau,\pi,\delta)$, and
stamps what its continuation brings to top level with $t+\len{\delta}$; a
chord stamps with $t+\max_i\len{\delta_i}$. Runs are onset-ordered: a set is
resolved only when its onset (the least onset of its candidates) is minimal
among all sets. The performance is the multiset of timed notes of positive
length.
\end{requirement}

\begin{requirement}[Engine neutrality]\label{req:neutral}
The engine \MUST{} know nothing of voices, fans, hands, servers, machines or
idioms. Those are definitions in the surface syntax that elaborate to terms.
(This is the property \texttt{skeinsim.py} states in its header, and it is what
makes Finding~8.3 --- the reflective voice plays the notes of the step
semantics --- a test of the calculus rather than of the code.)
\end{requirement}

\begin{remark}[What the note leaves to the implementation]
Three things are parameters of the semantics rather than parts of it: the chance
source (Definition~3.3 names it a parameter), the order in which simultaneously
due contention sets are resolved, and whether a set fires a maximal matching or
a single candidate. The last is \texttt{mode='one'} in the simulator and is
retained only as a diagnostic (T3). The first two are made explicit and recorded
in every trace (\S\ref{sec:chance}, \S\ref{ssec:sched}).
\end{remark}

%==============================================================================
\section{The textual format}\label{sec:syntax}

\subsection{Design position}

\begin{decision}[A dialect of the f1r3|@ng control flow sublanguage]\label{dec:dialect}
The surface syntax reuses the concrete syntax of the f1r3|@ng CFL wherever the
calculus has the same former: \texttt{for (y <- x where ...) \char`\{ P \char`\}},
\texttt{x!(...)}, \texttt{|}, \texttt{*x}, \texttt{@P}. What is new is the
decoration of names, the alphabet declarations, and a small elaboration-time
macro layer. The grammar is kept small enough to be restated later as a
\texttt{Theory Score} in the f1r3|@ng DDL once the Theory syntax lands in the
node; the parser crate is therefore a replaceable front end over a stable core
AST (\S\ref{sec:arch}).
\end{decision}

Files use the extension \texttt{.score} and are UTF-8. Comments are
\texttt{//} to end of line and \texttt{/* ... */}.

\subsection{Declarations}

A score file is a sequence of imports, declarations and definitions, followed by
exactly one \texttt{play} item giving the initial configuration.

\begin{lstlisting}[language=score]
score Twinkle

pitches   { r, C3, F3, G3, C4, D4, E4, F4, G4, A4 }   // r: the rest, required
durations { q = 1/4, h = 1/2 }                        // eps = 0 is implicit
timbres   { piano = gm(0) on 1, bass = gm(32) on 2 }  // GM program, MIDI channel
\end{lstlisting}

\begin{requirement}[Alphabets]\label{req:alph}
\begin{enumerate}[label=(\roman*),itemsep=1pt]
\item \texttt{pitches} lists $\Pit$ \emph{in order}; that order defines
  $\mathsf{step}(k)$ (\S\ref{ssec:atoms}). \texttt{r} \MUST{} be present and is
  unordered. A pitch written in scientific notation (\texttt{C4}, \texttt{F\#3},
  \texttt{Bb2}) gets its MIDI number implicitly (C4 = 60); any other identifier
  \MUST{} be given one, \texttt{deg7 = 67}.
\item \texttt{durations} gives each length as a non-negative rational literal
  (\texttt{1/6} for a triplet eighth, \texttt{3/8} for a dotted quarter).
  \texttt{eps = 0} is always present and \MUST{} not be redeclared. Other
  zero-length durations are permitted.
\item Ranges \MAY{} abbreviate alphabets:
  \texttt{pitches \char`\{ r, scale major C4 .. C7 \char`\}},
  \texttt{durations \char`\{ w=1, h=1/2, q=1/4, e=1/8, s=1/16 \char`\}}.
  Expansion is at elaboration time; the alphabet is always finite.
\item Identifiers in $\Pit$, $\Dur$ and the timbre set \MUST{} be pairwise
  distinct; the elaborator rejects overlap.
\end{enumerate}
\end{requirement}

\subsection{Core process syntax}

\begin{center}
\small
\begin{tabular}{@{}lll@{}}
\toprule
calculus & surface & note \\
\midrule
$0$ & \texttt{0} & \\
$\nm{P}{\tau}{d}$ & \texttt{<@P, tau, d>} & parenthesise compound \texttt{P} \\
$x!(Q,\tau,d)$ & \texttt{x!(Q, tau, d)} & \\
$\mathsf{for}(y\leftarrow x\;\mathsf{where}\;\psi)P$ &
  \texttt{for (y <- x where psi) \char`\{ P \char`\}} & clause optional \\
join & \texttt{for (y1 <- x1 \& y2 <- x2 where psi) \char`\{P\char`\}} & chord \\
$P\mid Q$ & \texttt{P | Q} & \\
$\ast x$ & \texttt{*x} & \\
stamp & \texttt{at 3/4 \char`\{ P \char`\}} & top level of \texttt{play} only \\
\bottomrule
\end{tabular}
\end{center}

\noindent The underscore \texttt{\_} is a binder that is not used. In the
\texttt{play} item every top-level receipt and message is stamped $0$ unless an
\texttt{at} says otherwise (an entrance).

\subsection{Definitions: the elaboration-time macro layer}

The note builds everything from abbreviations: $\mathsf N$, $\Rec$, $\Fan$,
$\mathsf{Step}$, $\Server$, $\Voice$, $\Call$, $\Resp$. The surface syntax needs
exactly this facility and no more.

\begin{lstlisting}[language=score]
def N(L: proc, tau: timbre, p: pitch, d: dur; K: proc) =
  for (_ <- <@L, tau, p>) { K } | <@L, tau, d>!(0, tau, r)
\end{lstlisting}

\begin{requirement}[Definitions are macros]\label{req:macros}
\begin{enumerate}[label=(\roman*),itemsep=1pt]
\item Parameters are sorted: \texttt{proc}, \texttt{name}, \texttt{timbre},
  \texttt{pitch}, \texttt{dur}, \texttt{clause}, and \texttt{cont} (a
  continuation: a definition of arity one in a name, passed by name).
\item Definitions are expanded at elaboration time and \MUST{} form an acyclic
  call graph. A recursive definition is rejected with a diagnostic pointing to
  \S4.2 of the note: unbounded behaviour is obtained by reflection, not by
  recursion (Remark~2.1).
\item Expansion is capture-avoiding. The elaborated term is a closed term of
  the core syntax and nothing else.
\end{enumerate}
\end{requirement}

\paragraph{Finite products.} $\prod_{v\in\Dur^+}\prod_{t\in\Pit}$ is written as
a comprehension that expands to a parallel composition:
\begin{lstlisting}[language=score]
def Fan(x: name, tau: timbre) =
  par v in durations+, t in pitches {
    <@*x, tau, v>!( <@*x, tau, v>!(*x, tau, t), tau, t )
  }
\end{lstlisting}
Generators range over \texttt{pitches}, \texttt{pitches-r} (the ordered
pitches), \texttt{durations}, \texttt{durations+} (positive length), and
explicit lists.

\paragraph{Lines.} A written part is a right-nested chain of written notes. The
derived form
\begin{lstlisting}[language=score]
line(Lm, piano) [ C4 q, C4 q, G4 q, G4 q, A4 q, A4 q, G4 h ]
\end{lstlisting}
elaborates to $\mathsf N_{Lm,\mathsf{piano}}(\mathrm{C4},\mathsf q;\,\cdots\,
\mathsf N_{Lm,\mathsf{piano}}(\mathrm{G4},\mathsf h;\,0)\cdots)$. An optional
trailing \texttt{then K} replaces the final $0$ with a continuation, which is
how a written phrase hands off to a responder (\S6.5 of the note).

\paragraph{Locations.} A \emph{base} is written \texttt{base "Lm"} and
elaborates to the closed receipt
$\mathsf{for}(\_\leftarrow\nm{0}{\tau_\bot}{\rest})\,0$ tagged with the string,
where $\tau_\bot$ is a reserved timbre no message may use. Distinct tags give
inequivalent bases.

\subsection{The standard library}

\begin{requirement}[\texttt{std.score}]\label{req:std}
The distribution \MUST{} ship \texttt{std.score}, written entirely in the
surface syntax, defining: \texttt{N}, \texttt{line}; \texttt{Rec}, \texttt{Fan},
\texttt{Seed} ($H_0$); \texttt{Step}, \texttt{Req}, \texttt{Server},
\texttt{Refresh}, \texttt{Voice} (\S4.2 of the note); \texttt{Call},
\texttt{Resp} (handoff); \texttt{HocketA}/\texttt{HocketB} (alternation); and
the dual (rhythm-leads) variants of \texttt{Fan} and \texttt{Step}.
\texttt{Voice} \MUST{} be the reflective voice, with no engine support: its
server keeps its own code on a code channel and re-sends it
(Proposition~4.4).
\end{requirement}

\begin{lstlisting}[language=score]
// std.score, excerpt: the reflective voice
def Req(z: name, tau: timbre, BS: proc) = <@BS, tau, eps>!(*z, tau, r)

def Server(tau: timbre, psi: clause, BS: proc, BC: proc) =
  for (w <- <@BS, tau, r>) {
    Step(w, tau, psi, Req(_, tau, BS)) | Refresh(tau, BC)
  }

def Refresh(tau: timbre, BC: proc) =
  for (c <- <@BC, tau, r>) { *c | <@BC, tau, eps>!(*c, tau, r) }

def Voice(tau: timbre, psi: clause, p0: pitch, u0: dur, tag: string) =
  Server(tau, psi, base tag ++ "/S", base tag ++ "/C")
  | <@(base tag ++ "/C"), tau, eps>!(Server(tau, psi, base tag ++ "/S",
                                            base tag ++ "/C"), tau, r)
  | <@(base tag ++ "/S"), tau, eps>!(Seed(base tag, tau, p0, u0), tau, r)
\end{lstlisting}

\subsection{Clauses and factors}

Clause syntax is given in \S\ref{sec:clauses}. At the top level a score may name
factors so that style, register and physicality are written once and combined:
\begin{lstlisting}[language=score]
factor stepwise = table step { -2..2: 3, _: 3/10 }
factor country  = machine pitch { C4 -> D4 @ 2, C4 -> E4 @ 3, D4 -> E4 @ 1, ... }
factor rhythm   = machine dur uniform

play Voice(piano, country * stepwise * rhythm, C4, q, "A")
   | Voice(bass,  walking * rhythm,            C3, h, "B")
\end{lstlisting}

\subsection{Twinkle, in full}

\begin{lstlisting}[language=score]
score Twinkle
import std
pitches   { r, C3, F3, G3, C4, D4, E4, F4, G4, A4 }
durations { q = 1/4, h = 1/2 }
timbres   { piano = gm(0) on 1, bass = gm(32) on 2 }

play line(base "Lm", piano) [ C4 q, C4 q, G4 q, G4 q, A4 q, A4 q, G4 h,
                              F4 q, F4 q, E4 q, E4 q, D4 q, D4 q, C4 h ]
   | line(base "Lb", bass)  [ C3 h, C3 h, F3 h, C3 h, F3 h, C3 h, G3 h, C3 h ]
\end{lstlisting}

\begin{requirement}[Golden score]\label{req:golden}
This file \MUST{} play, under every scheduler, the 22 timed notes of the table in
\S2.5 of the note, ending at 4 whole notes (conformance test T12).
\end{requirement}

\subsection{Diagnostics and lints}

The elaborator \MUST{} reject: overlapping alphabets; undeclared data or
timbres; recursive definitions; ill-sorted arguments; use of $\tau_\bot$;
clauses outside the checkable fragment (\S\ref{ssec:fragment}). It \SHOULD{}
warn when a message's payload timbre differs from its subject's timbre (such a
message can never communicate), and when a top-level receipt and message of the
same polarity share a location (they never communicate,
Proposition~2.4(ii)).

%==============================================================================
\section{Architecture and data model}\label{sec:arch}

\subsection{Crates}

\begin{center}
\small
\begin{tabular}{@{}l>{\raggedright\arraybackslash}p{9.2cm}@{}}
\toprule
crate & responsibility \\
\midrule
\texttt{score-core}   & data alphabets, hash-consed terms, congruence normal
  form, substitution, names and locations. No I/O. \\
\texttt{score-syntax} & lexer, parser, surface AST, macro elaboration,
  \texttt{std.score}, pretty-printer, diagnostics with spans. \\
\texttt{score-logic}  & crisp and graded formula ASTs, the checkable-fragment
  type, candidate views, evaluation, tables and machines. \\
\texttt{score-engine} & the soup, candidate generation, contention sets,
  maximal matchings, resolution, time, scheduling, collection, trace events. \\
\texttt{score-chance} & chance sources: PRNG digit streams, spigot digit
  streams, exact interval decoding, human choice, replay. \\
\texttt{score-render} & performance and trace serialisation, Standard MIDI
  File writer, live MIDI (feature \texttt{live}). \\
\texttt{f1r3score}    & the binary: command line, wiring, exit codes. \\
\texttt{score-conformance} & the ported skeinsim suite and golden files. \\
\bottomrule
\end{tabular}
\end{center}

\begin{requirement}[Portability]\label{req:port}
\texttt{score-core}, \texttt{score-syntax}, \texttt{score-logic},
\texttt{score-engine} and \texttt{score-chance} (without the spigot feature)
\MUST{} build for \texttt{wasm32-unknown-unknown} and \MUST{} perform no I/O,
so that the same engine can be embedded in a browser page or an Apple Vision
Pro client. They \SHOULD{} have no dependencies beyond \texttt{num-bigint} /
\texttt{num-rational}. The workspace targets stable Rust, edition 2021, with
\texttt{rust-version} pinned.
\end{requirement}

\subsection{Terms and hash-consing}

\begin{lstlisting}[language=Rust]
pub struct Datum(u16);                 // index into Pi + Delta (disjoint)
pub enum Polarity { Channel, CoChannel }
pub struct Timbre(u16);
pub struct ProcId(u32);                // hash-consed, congruence-normal process
pub struct Loc { pub quote: ProcId, pub timbre: Timbre }

pub enum Name { Var(VarIdx), Quote { proc: ProcId, timbre: Timbre, datum: Datum } }

pub enum Proc {                        // normal form only
    Nil,
    Par(Box<[ProcId]>),                // flattened, >= 2, sorted by id, no Nil
    Recv { binds: Box<[(Name /*subject*/)]>, clause: ClauseId, body: ProcId },
    Send { subj: Name, payload: ProcId, ptimbre: Timbre, carry: Datum },
    Drop(Name),                        // only on a variable: *<@P,..> normalises to P
}

pub struct Arena { /* Vec<Proc>, HashMap<Proc, ProcId> */ }
\end{lstlisting}

\begin{requirement}[Congruence by construction]\label{req:nf}
The arena \MUST{} intern only normal forms: parallel compositions flattened,
$0$ removed, children sorted by \texttt{ProcId} (a multiset), singletons
unwrapped; $\ast\nm{P}{\tau}{d}$ replaced by $P$; bound variables in de Bruijn
form so that $\alpha$-equivalent terms are identical. Then $P\equiv Q$ iff
their ids are equal, and location equivalence is integer equality. A property
test \MUST{} check this against an independent, unoptimised congruence decision
procedure on random terms.
\end{requirement}

\begin{remark}[Why sorting by id is sound]
Children are interned before their parent, so sorting by id is a canonical
order on already-normal children. Ids depend on interning order, which differs
between runs of different scores; they are never serialised as identities.
Serialised terms and digests use a structural encoding
(\S\ref{ssec:digest}).
\end{remark}

\begin{requirement}[Substitution]\label{req:subst}
Firing substitutes a closed name $\nm{Q}{\tau}{d}$ for the bound variable and
re-normalises; occurrences of $\ast y$ become $Q$ and names built from $y$ are
compared up to congruence thereafter. Substitution \MUST{} be memoised on
(body, name) since the same body is instantiated repeatedly by a voice.
\end{requirement}

\subsection{The soup}

The running configuration is not stored as a term. As in the simulator, it is a
\emph{soup}: the top-level receipts and messages, each an \emph{occurrence} with
a stable occurrence number, a stamp and (for receipts) its clause, indexed by
location.

\begin{lstlisting}[language=Rust]
pub struct RecvOcc { id: OccId, subjects: Box<[Name]>, clause: ClauseId,
                     body: ProcId, stamp: Time, stream: StreamKey }
pub struct SendOcc { id: OccId, subj: Name, payload: ProcId,
                     ptimbre: Timbre, carry: Datum, stamp: Time }
pub struct Soup {
    recv: BTreeMap<Loc, Vec<RecvOcc>>,
    send: BTreeMap<Loc, Vec<SendOcc>>,
    dirty: BTreeSet<Loc>,               // locations whose candidates may have changed
    next_occ: OccId,
}
\end{lstlisting}

\begin{requirement}[Determinism of iteration]\label{req:btree}
Every collection whose iteration order can influence a run \MUST{} be ordered
(\texttt{BTreeMap}, sorted \texttt{Vec}), keyed by occurrence number or by a
structural key. Occurrence numbers are assigned in firing order. Consequently a
score, a chance source and a scheduler determine the run exactly.
\end{requirement}

Putting a process into the soup (\emph{spawning} at stamp $t$) decomposes its
normal form: \texttt{Par} spawns each child; \texttt{Recv} and \texttt{Send}
become occurrences at their subjects' locations; \texttt{Nil} does nothing. A
top-level $\ast y$ cannot occur in a closed term.

\subsection{Digests}\label{ssec:digest}

Every elaborated score has a \emph{digest}: SHA-256 over a canonical,
id-independent serialisation of its term, alphabets and lengths. The digest
heads every trace and performance file and is the identity under which a score
would later be stored as a content-addressed value on a shard.

%==============================================================================
\section{The clause language}\label{sec:clauses}

\subsection{Grammar}

The two sorts of Paper~II, with the atoms of Definition~3.1 of the note.
\begin{center}
\small
\begin{tabular}{@{}lll@{}}
\toprule
sort & surface & meaning \\
\midrule
graded $\psi$ & \texttt{3/2}, \texttt{0.3} & constant $w\in\Qnn$ \\
 & \texttt{[beta]} & $\crisp\beta$ \\
 & \texttt{psi * psi}, \texttt{psi + psi} & $\otimes$, $\oplus$ \\
 & \texttt{<>psi} & graded behavioural modality $\langle K\rangle_\V$ \\
 & \texttt{mu X. psi}, \texttt{nu X. psi} & graded fixed points \\
 & \texttt{table key \char`\{ ... \char`\}}, \texttt{machine ...} & sugar, below \\
 & factor name & reference to a named factor \\
crisp $\beta$ & \texttt{true}, \texttt{!b}, \texttt{b \&\& b}, \texttt{b || b} &
  propositional \\
 & atoms & \S\ref{ssec:atoms} \\
 & \texttt{<>b}, \texttt{leads(b)}, \texttt{nu X. b} & behavioural \\
spatial $\sigma$ & \texttt{true}, \texttt{0}, \texttt{s | s}, \texttt{!s}, \texttt{s \&\& s} &
  over processes \\
 & \texttt{<n, tau, d>!(s, tau, d)} & $\mathsf{out}$: a message \\
 & \texttt{for(<n, tau, d>) s} & $\mathsf{in}$: a receipt \\
 & \texttt{@s} & name predicate \\
\bottomrule
\end{tabular}
\end{center}
\noindent In message and receipt patterns, \texttt{\_} matches any location,
timbre or datum; \texttt{n} is itself a spatial formula applied to the quoted
location, so patterns nest. Graded formulae have no negation (Paper~II,
Definition~2.4).

\subsection{Atoms and the candidate view}\label{ssec:atoms}

A clause is evaluated on a candidate $c$ against the configuration in which it
stands. The view is
\begin{lstlisting}[language=Rust]
pub struct View<'a> {
    pub timbre: Timbre, pub pitch: Datum, pub dur: Datum,   // the note nu(c)
    pub carry: Datum,                                       // kappa(c)
    pub loc: ProcId,                                        // @S
    pub passes: ProcId,                                     // Q of <@Q,tau,d>
    pub chord: &'a [ChordSlot],                             // joins only
    pub residue: Lazy<'a, Residue>,                         // behavioural atoms only
}
\end{lstlisting}

\begin{center}
\small
\begin{tabular}{@{}l>{\raggedright\arraybackslash}p{9.3cm}@{}}
\toprule
atom & holds when \\
\midrule
\texttt{pitch(s)}, \texttt{dur(v)}, \texttt{carry(e)} & the note's pitch, its
  duration, the carried datum equal the argument \\
\texttt{step(k)} & carry and pitch are ordered pitches and
  $\mathrm{ord}(\kappa)-\mathrm{ord}(\pi)=k$ \\
\texttt{at(s)} & the location's quote satisfies $\sigma$ \\
\texttt{passes(s)} & the passed process satisfies $\sigma$ \\
\texttt{held(t)} & \texttt{at(<\_,\_,\_>!(true, \_, t))} (derived) \\
\texttt{prev(u)} & \texttt{at(<\_,\_,u>!true)} (derived) \\
\texttt{back(j, e)} & the record $j$ levels into the past carries $e$ (derived) \\
\texttt{last[m1 .. mk]} & $\mathsf{at}(\mathsf{Last}(m))$, \S6.1 of the note (derived) \\
\texttt{pitch\#i}, \texttt{dur\#i}, \texttt{carry\#i} & per-pattern atoms of a
  chord receipt \\
\texttt{leads(b)} & the residue of firing $c$, in the present configuration,
  satisfies $\beta$ \\
\bottomrule
\end{tabular}
\end{center}

\begin{requirement}[Derived atoms are sugar]\label{req:derived}
\texttt{held}, \texttt{prev}, \texttt{back}, \texttt{last} \MUST{} elaborate to
spatial formulae. The evaluator \MAY{} compile them to direct walks over the
record chain (as the simulator's \texttt{back} and \texttt{repeat} do), but a
test \MUST{} check the compiled and the formula forms agree on random views,
exactly as T9 does for tables.
\end{requirement}

\subsection{Tables and machines}

\texttt{table key \char`\{ k: w, ..., \_: w \char`\}} is a sum over mutually exclusive crisp
conjunctions evaluated by lookup, keyed by \texttt{(pitch, carry)},
\texttt{(prev, dur)}, \texttt{step}, or any tuple of atoms that are functions of
the view. \texttt{\_} is the default (absent means 0). Ranges such as
\texttt{-2..2} are allowed on \texttt{step}.

\texttt{machine pitch \char`\{ s -> t @ w, ... \char`\}} and \texttt{machine dur \char`\{ u -> v @ w, ...\char`\}}
are Example~5.2's factors: the pitch machine is a table keyed by
\texttt{(pitch, carry)}, the duration machine a table keyed by
\texttt{(prev, dur)}. The arrow syntax is that of the \texttt{Machines} Theory in
the drawn-distributions design, and \texttt{machine pitch import
"file.machine"} reads a canonical machine exported by the instrument, mapping
its state indices onto the declared pitch alphabet by position (with the rest as
the last state, as on the instrument). \texttt{machine dur uniform} is the
complete uniform machine.

\subsection{Evaluation}

\begin{requirement}[Exact weights]\label{req:exact}
Clause values are exact non-negative rationals (\texttt{BigRational}, with a
fast path for \texttt{u64/u64} that promotes on overflow). Decimal literals are
read exactly (\texttt{0.3} is $3/10$). No floating point enters a weight, a
normalisation, or a choice.
\end{requirement}

\begin{requirement}[Purity]\label{req:pure}
Evaluation is a pure function of the view and the configuration. Evaluating a
clause on a candidate that does not fire changes nothing.
\end{requirement}

$\otimes$ short-circuits on $0$; $\oplus$ sums; $\crisp\beta$ is $1$ or $0$.
Graded fixed points are evaluated only inside the checkable fragment, by
iteration to the exact fixed point over a finite state space; if the iteration
does not reach an exact fixed point within the declared bound the score is
rejected at elaboration, never approximated at run time.

\subsection{The checkable fragment}\label{ssec:fragment}

The ideal logic defines; the fragment decides. As in \texttt{rho-weighted}, the
distinction is a type, not a comment.

\begin{lstlisting}[language=Rust]
pub struct Checked<F> { formula: F, class: Locality, depth: usize }
pub enum Locality { Local, Behavioural { window: usize } }
impl<F> Checked<F> { pub fn try_new(f: F, alph: &Alphabets) -> Result<Self, WhyNot>; }
\end{lstlisting}

\begin{requirement}[Fragment membership]\label{req:fragment}
A clause is admitted iff it is in one of two classes.
\begin{enumerate}[label=(\roman*),itemsep=1pt]
\item \textbf{Local}: no behavioural modality, no \texttt{leads}. Its value
  depends only on the view. Spatial formulae are decidable on the finite
  quoted terms involved; separating conjunction is decided by enumerating splits
  of a multiset and \MUST{} be bounded by a configured maximum split count.
\item \textbf{Window-behavioural}: behavioural modalities and fixed points are
  admitted when every spatial formula in \emph{every} clause of the score reads
  the past to depth at most $d$ (the depth grading of the rho-life notes) and
  the alphabets are finite. Then the state relevant to the formula is finite
  (\S6.2 of the note), and $\nu$ is computed exactly by Knaster--Tarski
  iteration over the abstract graph whose nodes are configurations with every
  location truncated to depth $d$ and stamps erased.
\end{enumerate}
Anything else is rejected with \texttt{WhyNot} naming the offending
subformula.
\end{requirement}

\begin{decision}[Behavioural modalities read the actual company]\label{dec:company}
In Proposition~6.3 the contexts $K$ range over ``the one-hole contexts supplied
by the rest of the score''. A player has exactly one such context at each
moment: the rest of the current configuration. The player evaluates
$\langle K\rangle$ and \texttt{leads} against that, and the \texttt{<>} surface
form carries no explicit context. Quantifying over hypothetical companies is an
analysis task, not a playing task, and is out of scope.
\end{decision}

\begin{requirement}[Locality classes drive caching]\label{req:locality}
Local clauses depend only on the candidate, so a location's candidate list is
recomputed only when that location is dirty. If any receipt in the soup carries
a behavioural clause, every such receipt's candidates are re-evaluated after
every resolution. The engine \MUST{} report in the trace header which regime
the run used.
\end{requirement}

%==============================================================================
\section{The engine}\label{sec:engine}

\subsection{One resolution}

\begin{enumerate}[leftmargin=*,itemsep=2pt]
\item \textbf{Candidates.} For each dirty location, pair every receipt with
  every message of opposite polarity whose payload timbre equals the location's
  timbre (for a join, every injective assignment of messages to its patterns,
  all at that location). Build the view, evaluate the clause, keep the
  candidates of positive weight.
\item \textbf{Contention sets.} Union--find over shared receipt or message
  occurrence numbers, per location. Since communication requires equal locations,
  a contention set never spans two locations; the engine \MUST{} assert this.
\item \textbf{Onsets.} A candidate's onset is $\max(t_r,t_m)$ (for a join, the
  maximum over all its stamps); a set's onset is the least of its candidates'.
\item \textbf{Due sets.} Let $t_{\min}$ be the least onset over all sets. The
  due sets are those with onset $t_{\min}$; the scheduler picks one
  (\S\ref{ssec:sched}).
\item \textbf{Alternatives.} Enumerate the maximal matchings of the chosen set
  in canonical order, each with weight $\prod_{c\in\mathcal M}\psi\llbracket
  c\rrbracket$.
\item \textbf{Choice.} Ask the set's chance source for an index with
  probability proportional to the weights (\S\ref{sec:chance}).
\item \textbf{Firing.} For each candidate in the matching: remove both
  occurrences, emit the timed note, form the substituted continuation and spawn
  it at the new stamp. Mark affected locations dirty. Append one trace event.
\end{enumerate}

\begin{requirement}[Canonical order of alternatives]\label{req:canon}
Candidates are ordered by (receipt occurrence number, message occurrence
numbers); matchings by the lexicographic order of their sorted candidate lists.
Together with R\ref{req:btree}, this makes ``the $i$-th alternative'' a
well-defined, replayable quantity.
\end{requirement}

\subsection{Maximal matchings}

\begin{requirement}[Stars and the general case]\label{req:matchings}
\begin{enumerate}[label=(\roman*),itemsep=1pt]
\item \textbf{Star fast path.} If every candidate in a set shares one receipt,
  or every candidate shares one message, the maximal matchings are the single
  candidates. This covers every voice, every chord receipt and every null note,
  and \MUST{} be $O(\text{set size})$ with no enumeration.
\item \textbf{General case.} Otherwise enumerate maximal matchings of the
  bipartite (for joins, hypergraph) contention structure by backtracking with
  maximality pruning. The number of alternatives \MUST{} be bounded by a
  configurable limit (default $10^5$); exceeding it is a run-time error naming
  the location, never a silent truncation.
\item \textbf{Mode \texttt{one}.} A diagnostic mode, off by default, in which
  a set fires a single candidate chosen with probability proportional to its
  weight. It exists to reproduce T3's second column and is recorded in the
  trace header.
\end{enumerate}
\end{requirement}

\subsection{Scheduling simultaneous sets}\label{ssec:sched}

When several sets are due at the same onset, the note is indifferent for
independent parts (Theorem~2.6, Theorem~7.1) and the order matters only where
parts contend. The player makes the order an explicit, recorded parameter.

\begin{decision}[Schedulers]\label{dec:sched}
\texttt{canonical} (default): the due set whose location has the least
structural key (digest of the quote, then timbre). \texttt{random:SEED}: a
dedicated PRNG, separate from every chance source, shuffles the due sets.
\texttt{by-timbre:A,B,...}: a stated priority, for the Twinkle-style ``melody
first / bass first'' tests. The scheduler never consumes digits from a musical
chance source.
\end{decision}

\subsection{Time}

Stamps are exact rationals. Continuations spawned by a firing receive the stamp
of R\ref{req:time}. Joins use the ``longest duration'' convention of
Definition~5.5; since the note lists this as an open thread, the convention is a
single function in the engine, not a scattered assumption.

\subsection{Termination and productivity}

A run stops at the first of: no available candidate anywhere (the score has
ended or every part has fallen silent); a \texttt{--notes N} limit on notes of
positive length; a \texttt{--until T} limit on onsets (whole notes); a
\texttt{--max-steps} limit on resolutions.

\begin{requirement}[Unproductive loops]\label{req:productive}
The note does not yet give a syntactic productivity check (open thread 6). The
engine \MUST{} therefore bound the number of consecutive zero-length notes at one
onset (default $10^5$) and stop with an \texttt{unproductive} error naming the
locations involved. A voice plays exactly two null notes per note
(Proposition~4.4), so the bound is never approached by a well-formed voice.
\end{requirement}

\subsection{Collection of inert locations}

A voice leaves $|\Dur^+|\cdot|\Pit|-1$ inert messages and $|\Pit|-1$ inert
hands behind at every note (Findings~8.1 and 8.3). Keeping them is correct but
costs memory linear in the length of the piece.

\begin{requirement}[Freshness collection]\label{req:gc}
Collection is off by default. With \texttt{--gc freshness}, after a resolution
at a location whose quote is a message (a record, not a base) and which then has
no available candidate, the location's occurrences are dropped and its key is
added to a tombstone set. If anything is later spawned at a tombstoned location,
the engine \MUST{} stop with a \texttt{freshness-violated} error: the collection
was unsound for this score. For scores built from \texttt{std.score} voices this
cannot happen (Lemma~4.5).
\end{requirement}

\begin{remark}[Why not reachability]
Collecting a location because no live term mentions it is not sound in this
calculus: any process can quote any closed term it can build, and bases are
built from nothing. Freshness is a theorem about how voices use locations, so
the player relies on it only where it is checked.
\end{remark}

%==============================================================================
\section{Chance sources}\label{sec:chance}

\subsection{One interface}

\begin{lstlisting}[language=Rust]
pub trait Chance {
    /// Choose an index into `weights` (exact, all > 0) with probability
    /// proportional to weight. Returns the index and what was consumed.
    fn choose(&mut self, weights: &[Weight], ctx: &ChoiceCtx) -> (usize, Consumed);
    fn describe(&self) -> SourceSpec;          // recorded in the trace header
}
pub trait DigitStream { fn base(&self) -> u32; fn next(&mut self) -> u32; }
\end{lstlisting}

\begin{decision}[Every automatic source is a digit stream]\label{dec:digits}
Every non-human source is a stream of digits in some base, and every choice is
made by the exact interval decoding of the drawn-distributions design
(Algorithm~1): cumulative weights are cleared to a common denominator, digits
narrow a $b$-adic interval, and the choice is the bin that contains it. There is
no floating point and no rejection sampling, so a given digit stream yields the
same choices on every machine. $K_{\max}$ (default 64 digits per choice) and its
deterministic fallback are part of the source's specification.
\end{decision}

\begin{center}
\small
\begin{tabular}{@{}l>{\raggedright\arraybackslash}p{9.2cm}@{}}
\toprule
spec & source \\
\midrule
\texttt{prng:SEED} & a specified counter-based generator (SplitMix64 feeding
  xoshiro256**, fixed in this specification so it can never drift with a
  dependency), emitting base-$2^{32}$ digits \\
\texttt{spigot:C/B} & digits of the fractional part of constant \texttt{C}
  $\in\{$\texttt{pi, e, ln2, liouville, champernowne, thue-morse}$\}$ in base
  \texttt{B}, from \texttt{spigot\_stream} (feature \texttt{spigot}); the
  integer part is skipped \\
\texttt{human} & an interactive chooser (\S\ref{ssec:human}) \\
\texttt{replay:FILE} & the choice indices recorded in a trace \\
\bottomrule
\end{tabular}
\end{center}

\subsection{Which source a choice uses}

Theorem~7.1(iii) holds \emph{surely} only when each voice has its own source
(T2 uses one stream per voice). The player routes choices through \emph{stream
keys}.

\begin{decision}[Stream keys]\label{dec:streams}
Every receipt occurrence carries a stream key, inherited from the receipt that
spawned it. The key is the receipt's timbre by default, and \MAY{} be set
explicitly with an annotation \texttt{for \#A (...)} or on a definition
application \texttt{Voice(...) \#A}; annotations are metadata and do not change
the elaborated term's congruence class or digest. The command line maps keys to
sources: \texttt{--chance prng:7 --stream A=spigot:pi/16 --stream B=prng:12}.
A contention set uses the stream of its least receipt occurrence. Sets with a
single alternative (every null note) consume nothing.
\end{decision}

\subsection{Factorised resolution}

Lemma~5.3 (factorised resolution) lets a star whose weights factor as $f(a)g(b)$
be resolved by two independent draws. T4 depends on this: pitch from the
hexadecimal digits of $\pi$, duration from the base-5 digits of $e$, one digit
each.

\begin{requirement}[Factor routing]\label{req:factor}
When a star's clause is a product of factors that are functions of disjoint
parts of the view (the elaborator can see this syntactically: a
\texttt{(pitch, carry)} machine times a \texttt{(prev, dur)} machine), the
engine \MAY{} resolve it factor by factor, each factor drawing from its own
stream (\texttt{--stream A.pitch=spigot:pi/16 --stream A.dur=spigot:e/5}). This
is an equivalent resolution by Lemma~5.3 and \MUST{} be recorded in the trace.
Without a per-factor stream the joint draw is used.
\end{requirement}

\subsection{Human choice}\label{ssec:human}

Proposition~6.4 makes the improviser a chance source that ignores the
probabilities and chooses among available candidates.

\begin{requirement}[Human chooser]\label{req:human}
The \texttt{human} source presents, for each set with more than one
alternative, the alternatives as notes (pitch, duration, carried pitch), with
their normalised probabilities shown for guidance, and accepts an index. Only
alternatives of positive weight are offered, so every trace a human produces is
admissible, and with a viability factor never runs out
(Proposition~6.4(iii)). Mixed performance --- a human on one stream key, an
algorithmic source on the others --- is the default when \texttt{human} is bound
to a single key. In this version the human chooser is step-wise, not real time.
\end{requirement}

%==============================================================================
\section{Outputs}\label{sec:outputs}

\subsection{Performance}

The performance is printed as a table by default and as JSON with
\texttt{--out perf.json}:
\begin{lstlisting}
{ "score": "sha256:...", "end": "4",
  "notes": [ { "onset": "0", "timbre": "piano", "pitch": "C4", "dur": "q", "len": "1/4" },
             { "onset": "0", "timbre": "bass",  "pitch": "C3", "dur": "h", "len": "1/2" }, ... ] }
\end{lstlisting}
Notes are sorted by (onset, timbre, pitch, duration) so that two runs with the
same performance produce byte-identical files; this is how interleaving
invariance is tested. Null notes are excluded (they are not part of the
performance); \texttt{--include-null} adds them for diagnosis.

\subsection{Trace}

\begin{requirement}[Replayable trace]\label{req:trace}
With \texttt{--trace run.jsonl} the player writes JSON Lines: a header
(format version, player version, score digest, alphabets, every source spec,
stream bindings, scheduler, matching mode, locality regime, collection mode),
then one event per resolution: step number, onset, location key, set size,
number of alternatives, chosen index, digits consumed per stream, and the notes
emitted. \texttt{f1r3score replay} \MUST{} reproduce the performance and the
event sequence from the score and the trace alone, and \MUST{} fail if the
score's digest differs from the header's.
\end{requirement}

\subsection{MIDI}

\begin{requirement}[Standard MIDI File]\label{req:midi}
\texttt{--midi out.mid} writes a format-1 file with one track per timbre, a
program change from the timbre's declaration, and note-on/note-off pairs at
$\mathrm{ticks}(t)$ and $\mathrm{ticks}(t+\len\delta)$. Rests produce no
events. The division (ticks per quarter) is the least value making every onset
and length in the performance an integer number of ticks, if that is at most
32767; otherwise 960 with rounding and a warning. Tempo is \texttt{--bpm}
(default 120), a rendering parameter only.
\end{requirement}

\begin{finding}[\texttt{spigot\_midi} cannot be reused as is]
\texttt{spigot\_midi::MidiTrack} writes format 0 from a sequence of
(pitch, duration, velocity) notes with no onsets: it is monophonic and
sequential by construction. A score is polyphonic, with explicit onsets and
simultaneous strikes. The player therefore needs its own small writer. It
\SHOULD{} follow \texttt{spigot\_midi}'s dependency-free style, reuse its
\texttt{GeneralMidi} enumeration for program names, and \MAY{} later be
upstreamed to replace it.
\end{finding}

\subsection{Live MIDI}

With feature \texttt{live} and \texttt{--live [PORT]}, a playback thread
receives timed notes from the engine through a bounded channel and sends them to
a MIDI port via \texttt{midir}, as \texttt{leap\_spigot::player} does. The
engine runs ahead of the clock by a configurable look-ahead in whole notes. The
engine never waits on the clock: live output is a rendering of the same run,
and its trace is identical to an offline run's.

%==============================================================================
\section{Command line}\label{sec:cli}

\begin{lstlisting}
f1r3score check   piece.score              # parse, elaborate, lint, report fragment and digest
f1r3score expand  piece.score              # print the elaborated core term
f1r3score play    piece.score [options]    # run once
f1r3score replay  piece.score run.jsonl [--midi out.mid]
f1r3score machine import skein.machine --pitches piece.score   # print as a factor

options for play:
  --chance SPEC            default source (prng:SEED | spigot:C/B | human)
  --stream KEY=SPEC        bind a stream key (repeatable; KEY.pitch / KEY.dur for factors)
  --scheduler S            canonical | random:SEED | by-timbre:A,B,...
  --notes N | --until T | --max-steps N
  --mode max|one           matching mode (default max)
  --gc off|freshness       collection (default off)
  --max-alternatives N     bound on maximal matchings per set (default 100000)
  --out FILE  --trace FILE  --midi FILE  --bpm N  --live [PORT]  --include-null
\end{lstlisting}

\noindent Exit codes: 0 success; 1 usage; 2 parse or elaboration error; 3
run-time error (\texttt{unproductive}, \texttt{freshness-violated}, alternative
limit); 4 replay mismatch.

%==============================================================================
\section{Conformance}\label{sec:conformance}

\subsection{The ported suite}

Each skeinsim test becomes a Rust integration test in
\texttt{score-conformance}, reading its score from a \texttt{.score} file (not
constructed in Rust), so the suite tests the front end as well as the engine.
Statistical thresholds are those the note reports; tests are seeded and
deterministic.

\begin{center}
\small
\begin{tabular}{@{}l>{\raggedright\arraybackslash}p{5.2cm}>{\raggedright\arraybackslash}p{5.3cm}@{}}
\toprule
test & what & pass criterion \\
\midrule
T0 & inertness of the unplayed & one live location per step; exactly 79 stale
  messages per note (16 pitches, 5 durations) \\
T1 & law of one voice (pitch leads) & pitch, duration and joint deviations
  within sampling error: every $|z|\le 3.5$; no off-machine transition \\
T2 & harmony by independence & five schedulers, byte-identical performances \\
T3 & two-sided fitting & \texttt{max}: first resolution emits $\min(m,n)$ notes;
  \texttt{one}: 1, then $\min(m,n)-1$ more \\
T4 & instrument as a score & pitch = hex digits of $\pi$, duration = base-5
  digits of $e$, one digit each, 400 notes \\
T5 & law of the dual (rhythm leads) & as T1, via \textsc{Comm}$_2$ \\
T6 & simultaneity by contention & two notes per resolution; equal durations
  never \\
T7 & style $\otimes$ physicality & transitions match normalised product;
  $|z|\le 3.5$; no forbidden transition \\
T8 & idioms from the past & zero triple repeats under the crisp factor; no
  silence; motif completion within $0.03$ of analytic \\
T9 & table = formula & zero difference over 3000 random views (exact
  rationals) \\
T10 & reflective = step semantics & identical performance over 3000 notes;
  exactly two null notes per note, all $(\rest,\nulld)$ \\
T11 & reflective voice, inertness & at most one record location with a
  candidate; $|\Pit|-1$ dead hands per note \\
T12 & Twinkle & 22 notes, end 4, identical under both orders \\
\bottomrule
\end{tabular}
\end{center}

\noindent T10 is the central test. In the simulator it compares the reflective
voice against a Python-constant voice. Here both sides are scores: the
reflective \texttt{Voice} from \texttt{std.score} against an unrolled voice of
$k$ written steps, with the same clause and the same musical stream.

\subsection{Additional tests}

\begin{itemize}[leftmargin=*,itemsep=2pt]
\item \textbf{Congruence.} Normal-form interning agrees with an independent
  decision procedure on $10^4$ random term pairs, including $\ast@$ and
  $\alpha$-variants.
\item \textbf{Cross-implementation differential.} For T0--T12 scores and
  random small scores, the Rust player and \texttt{skeinsim.py}, driven by the
  same recorded choice indices, emit the same note sequence. The Python side
  gains a small adapter that accepts a choice list; this is the faithfulness
  obligation, kept as a permanent CI job.
\item \textbf{Replay.} Every conformance run's trace replays to an identical
  performance.
\item \textbf{Enforcement.} A window-behavioural score using
  $\mathsf{Viable}_\beta$ (Proposition~6.3) runs $10^4$ notes without violating
  $\beta$ and without falling silent; the same score without the viability
  factor is shown, on a fixture, to fall silent.
\item \textbf{Chords.} A chord receipt emits $k$ notes at one onset, and its
  continuation is stamped at the longest duration.
\item \textbf{Freshness guard.} A deliberately non-fresh score run with
  \texttt{--gc freshness} fails with \texttt{freshness-violated}.
\item \textbf{MIDI.} Twinkle's file is byte-identical to a golden file.
\end{itemize}

%==============================================================================
\section{Phasing}\label{sec:phases}

\begin{center}
\small
\begin{tabular}{@{}l>{\raggedright\arraybackslash}p{5.4cm}>{\raggedright\arraybackslash}p{5.2cm}@{}}
\toprule
phase & content & exit criteria \\
\midrule
P0 & workspace; \texttt{score-core} arena, normal forms, substitution, digests;
  parser for the core syntax and alphabets; pretty-printer &
  round-trip parse/print; congruence property test; Twinkle parses \\
P1 & engine with clause-free receipts: soup, candidates, contention, maximal
  matchings, time, schedulers; performance output; trace &
  T3, T6, T12 \\
P2 & local clauses: atoms, spatial formulae, tables, machines, exact weights;
  PRNG source with interval decoding; macros and \texttt{par} comprehensions &
  T0, T1, T5, T7, T8, T9 \\
P3 & \texttt{std.score}: reflective voice, handoff, hocket; freshness collection;
  stream keys & T2, T10, T11; freshness guard \\
P4 & MIDI writer; spigot sources and factor routing; replay; live MIDI &
  T4; replay; MIDI golden \\
P5 & window-behavioural fragment: \texttt{<>}, \texttt{leads}, $\nu$; chord
  receipts & enforcement and chord tests \\
P6 & human chooser; wasm build of the core crates; cross-implementation CI &
  differential job green; wasm build in CI \\
\bottomrule
\end{tabular}
\end{center}

\noindent P0--P3 deliver a player for every example in the note except
Proposition~6.3 and Definition~5.5; P4 makes it musically usable; P5 completes
the calculus as the note defines it.

%==============================================================================
\section{Decisions to confirm}\label{sec:open}

\begin{enumerate}[leftmargin=*,itemsep=3pt]
\item \textbf{Home of the code.} Candidates: a new repository
  \texttt{F1R3FLY-io/f1r3score}; a directory beside the note in
  \texttt{publications/f1r3score}; or \texttt{F1R3Games}. Since the calculus is
  deliberately not F1R3Skein, the proposal is a new repository, with the note's
  directory linking to it.
\item \textbf{Surface syntax.} Confirm Decision~\ref{dec:dialect}: a CFL
  dialect now, restated as \texttt{Theory Score} in the DDL later. In
  particular confirm the name syntax \texttt{<@P, tau, d>} and the file
  extension \texttt{.score}.
\item \textbf{Exact rationals by default.} Decision~\ref{dec:digits} and
  R\ref{req:exact} put bit-for-bit reproducibility ahead of speed. An
  \texttt{f64} mode could be added for large ensembles; the proposal is not to
  add one until a measured need appears.
\item \textbf{Stream keys.} Decision~\ref{dec:streams} defaults to timbre, with
  optional labels. Two voices sharing a timbre (hocket) would otherwise share a
  stream; confirm that labels are the right remedy rather than, say, keying by
  base.
\item \textbf{Scheduler default.} \texttt{canonical} makes runs reproducible
  without a seed. Confirm, or prefer a seeded random default.
\item \textbf{Behavioural modalities.} Decision~\ref{dec:company} evaluates them
  against the actual company, and R\ref{req:fragment} admits them only in the
  window fragment. Confirm that a bounded-lookahead approximation, which would
  not guarantee liveness, should stay out.
\item \textbf{Chord continuation.} The longest-duration convention of
  Definition~5.5 is isolated in one function; confirm it as the default while
  open thread 5 is pursued.
\item \textbf{Handoff record.} Proposition~6.5 notes that the outermost record
  of a handoff carries a pitch that was never played. Confirm that
  \texttt{last[...]} should not skip it automatically (the proposal: expose it,
  and give \texttt{std.score} a \texttt{call[...]} atom that skips it).
\item \textbf{Pitch naming.} Scientific pitch names with C4 = 60, and explicit
  MIDI numbers otherwise. Confirm, including enharmonic spellings as distinct
  data (so \texttt{F\#4} and \texttt{Gb4} are different pitches with the same
  MIDI number).
\end{enumerate}

\end{document}
